% JEDYNY PLIK, KTÓRY TRZEBA EDYTOWAĆ PRZED ROZPOCZĘCIEM PISANIA
% Język: polish lub english. Typ: engineering, master lub doctoral.
\newcommand{\thesislanguage}{polish}
\newcommand{\thesistype}{doctoral}
% Etap pracy: draft ostrzega, final zatrzymuje kompilacje przy brakach metadanych.
\newcommand{\thesisstatus}{draft}
% Sposob publikacji: digital (PDF) albo print (druk dwustronny).
\newcommand{\thesisprintmode}{digital}
% Wciecie pierwszego akapitu po naglowku: true albo false.
\newcommand{\indentfirstparagraph}{false}
% Opcjonalne wykazy po spisie tresci: true (wlaczony) albo false (wylaczony).
\newcommand{\showlistoffigures}{false}
\newcommand{\showlistoftables}{false}
\newcommand{\showlistoflistings}{false}
\newcommand{\showabbreviationlist}{false}
\newcommand{\showsymbollist}{false}
% Opcjonalne strony i oswiadczenia: true (wlaczone) albo false (wylaczone).
\newcommand{\showdedication}{false}
\newcommand{\showaideclaration}{false}
\newcommand{\showacknowledgements}{false}
\newcommand{\showfundingstatement}{false}
\newcommand{\showavailabilitystatement}{false}

\newcommand{\universitypl}{Politechnika Opolska}
\newcommand{\universityen}{Opole University of Technology}
\newcommand{\facultypl}{Wydział Elektrotechniki, Automatyki i Informatyki}
\newcommand{\facultyen}{Faculty of Electrical Engineering, Automatic Control and Informatics}
\newcommand{\titlepl}{Tytuł pracy w języku polskim}
\newcommand{\titleen}{Thesis title in English}
\newcommand{\authorname}{Imię i nazwisko}
\newcommand{\albumid}{000000}
\newcommand{\programmepl}{Nazwa kierunku}
\newcommand{\programmeen}{Study programme}
\newcommand{\supervisorname}{dr hab. inż. Imię Nazwisko, prof. uczelni}
\newcommand{\auxsupervisorname}{}
\newcommand{\city}{Opole}
\newcommand{\submissionyear}{2026}
\newcommand{\keywordspl}{słowo kluczowe 1, słowo kluczowe 2, słowo kluczowe 3}
\newcommand{\keywordsen}{keyword 1, keyword 2, keyword 3}
% Logo i jego szerokość na stronie tytułowej.
\newcommand{\universitylogo}{assets/branding/logo-weaii.jpeg}
\newcommand{\universitylogowidth}{0.58\textwidth}
